\documentclass{article}
\usepackage[utf8]{inputenc}
\usepackage{a4wide}

\begin{document}

\section*{สำรวจกล่องลูกอมสี}

คุณแม่ให้ลูก ๆ จำนวน N คนเอากล่องลูกอมมาคนละกล่อง\\

แต่ละกล่องมีลูกอมจำนวนหนึ่ง (เป็นจำนวนเต็ม)



คุณแม่อยากรู้ว่า ในกล่องทั้งหมดนี้มีกล่องที่มีจำนวนลูกอมเป็นเลขคู่กี่กล่อง



คำสั่ง:



รับจำนวนเต็ม N และจำนวนลูกอมในแต่ละกล่อง



นับและแสดงจำนวนกล่องที่มีจำนวนลูกอมเป็นเลขคู่



\subsection*{Input}
\begin{itemize}

\item

  บรรทัดแรกเป็นจำนวนเต็มบวก\texttt{N}(1 ≤ N ≤ 100) แทนจำนวนกล่องลูกอม

\item

  บรรทัดที่สองเป็นจำนวนเต็ม N ตัว โดยแต่ละตัวแทนจำนวนลูกอมในแต่ละกล่อง

\item

  จำนวนลูกอมในแต่ละกล่องจะอยู่ในช่วง 0 ถึง 1,000

\end{itemize}



\subsection*{Output}
แสดงจำนวนกล่องที่มีจำนวนลูกอมเป็นเลขคู่ ตามรูปแบบดังนี้:



\texttt{จำนวนกล่องที่มีลูกอมเลขคู่:\ X}



โดยที่\texttt{X}คือจำนวนกล่องที่จำนวนลูกอมเป็นเลขคู่



\end{document}
